% ============================================================
\chapter{NoSQL Databases}
\label{ch:nosql}
% ============================================================

\begin{intuition}[title=\textbf{Central idea}]
NoSQL (\emph{Not Only SQL}) databases abandon the strict relational model
to gain horizontal scalability, schema flexibility, and performance on
specific use cases. Four major families exist: key-value, document,
column-family, and graph.
\end{intuition}

% ------------------------------------------------------------
\section{Motivation and context}
% ------------------------------------------------------------

Relational databases reach their limits when facing:
\begin{itemize}
  \item massive volumes (petabytes),
  \item semi-structured or unstructured data,
  \item high availability and geographic distribution requirements,
  \item evolving schemas (agile development).
\end{itemize}

\begin{remark}
NoSQL does not mean ``anti-SQL'' but ``Not Only SQL.'' Many NoSQL databases
now offer a declarative query language (CQL for Cassandra, MQL for MongoDB,
Cypher for Neo4j).
\end{remark}

% ------------------------------------------------------------
\section{The CAP theorem}
% ------------------------------------------------------------

\begin{theorem}[CAP theorem (Brewer, 2000)]
\index{CAP theorem}
A distributed system can simultaneously guarantee at most two of the
following three properties:
\begin{enumerate}
  \item \textbf{Consistency} --- every read returns the most recent write.
  \item \textbf{Availability} --- every request receives a response (not
        necessarily the most recent data).
  \item \textbf{Partition tolerance} --- the system continues to operate
        despite network partitions.
\end{enumerate}
\end{theorem}

\begin{attention}
In practice, network partitions are unavoidable in a distributed system.
The real choice is therefore between \textbf{CP} (strong consistency,
e.g., HBase, MongoDB with \texttt{readConcern: majority}) and \textbf{AP}
(availability, e.g., Cassandra, DynamoDB in eventual mode).
\end{attention}

\begin{definition}[Eventual consistency]
\index{eventual consistency}
A system with \textbf{eventual consistency} guarantees that, in the absence
of new writes, all replicas converge to the same value after a finite delay.
\end{definition}

% ------------------------------------------------------------
\section{Key-value stores}
% ------------------------------------------------------------

\begin{definition}[Key-value store]
\index{key-value store}
A \textbf{key-value store} stores pairs $(k, v)$ where $k$ is a unique key
and $v$ is an opaque value (string, blob, serialized object). Access is by
key only: \texttt{GET}, \texttt{SET}, \texttt{DELETE}.
\end{definition}

\begin{example}
Redis: in-memory cache and data structure server:
\begin{minted}{python}
import redis

r = redis.Redis(host='localhost', port=6379, db=0)
r.set('user:1001:name', 'Alice')
r.set('user:1001:email', 'alice@example.com')
r.expire('user:1001:name', 3600)  # 1-hour TTL

name = r.get('user:1001:name')  # b'Alice'

# Advanced data structures
r.hset('user:1002', mapping={'name': 'Bob', 'age': '25'})
r.lpush('queue:jobs', 'task_42')
r.zadd('leaderboard', {'Alice': 1500, 'Bob': 1200})
\end{minted}
\end{example}

\begin{remark}
Typical use cases: application cache, user sessions, real-time counters,
message queues.
\end{remark}

% ------------------------------------------------------------
\section{Document stores}
% ------------------------------------------------------------

\begin{definition}[Document store]
\index{document store}
A \textbf{document store} stores semi-structured documents (JSON, BSON)
grouped into \emph{collections}. Each document has a flexible schema: two
documents in the same collection may have different fields.
\end{definition}

\begin{example}
MongoDB:
\begin{minted}{python}
from pymongo import MongoClient

client = MongoClient('mongodb://localhost:27017/')
db = client['university']
col = db['students']

# Insertion
col.insert_one({
    'last_name': 'Smith',
    'first_name': 'Alice',
    'grades': [
        {'course': 'Databases', 'grade': 92},
        {'course': 'Algorithms', 'grade': 88}
    ],
    'address': {'city': 'New York', 'zip': '10001'}
})

# Query with filter and projection
results = col.find(
    {'grades.grade': {'$gte': 90}},
    {'last_name': 1, 'first_name': 1, '_id': 0}
)

# Aggregation (equivalent to GROUP BY)
pipeline = [
    {'$unwind': '$grades'},
    {'$group': {'_id': '$last_name', 'avg_grade': {'$avg': '$grades.grade'}}}
]
col.aggregate(pipeline)
\end{minted}
\end{example}

\begin{remark}
MongoDB offers secondary indexes, multi-document transactions (since
version 4.0), and a rich aggregation language. It suits data with variable
schemas (product catalogs, CMS, IoT).
\end{remark}

% ------------------------------------------------------------
\section{Column-family stores}
% ------------------------------------------------------------

\begin{definition}[Column-family store]
\index{column-family store}
A \textbf{column-family store} organizes data by column families rather
than rows. Each row is identified by a partition key and can have a
variable number of columns.
\end{definition}

\begin{example}
Apache Cassandra:
\begin{minted}{python}
-- CQL (Cassandra Query Language)
CREATE KEYSPACE university WITH replication = {
    'class': 'SimpleStrategy', 'replication_factor': 3
};

CREATE TABLE university.grades (
    student_id UUID,
    course TEXT,
    semester INT,
    grade FLOAT,
    PRIMARY KEY ((student_id), course, semester)
) WITH CLUSTERING ORDER BY (course ASC, semester DESC);

-- The partition key (student_id) determines the storage node
-- Clustering columns order data within the partition
\end{minted}
\end{example}

\begin{remark}
Cassandra excels at massive distributed writes (logs, time series, IoT).
Modeling is query-driven: tables are designed based on the \texttt{SELECT}
queries you plan to execute.
\end{remark}

% ------------------------------------------------------------
\section{Graph databases}
% ------------------------------------------------------------

\begin{definition}[Graph database]
\index{graph database}
A \textbf{graph database} stores \textbf{nodes} (entities) and
\textbf{edges} (relationships) as first-class citizens. Graph traversals
are native operations running in $O(1)$ per relationship (independent of
graph size).
\end{definition}

\begin{example}
Neo4j and the Cypher language:
\begin{minted}{python}
// Create nodes and relationships
CREATE (a:Person {name: 'Alice', age: 30})
CREATE (b:Person {name: 'Bob', age: 25})
CREATE (a)-[:FRIENDS_WITH {since: 2020}]->(b)
CREATE (a)-[:WORKS_AT]->(:Company {name: 'TechCorp'})

// Find friends of friends (distance 2)
MATCH (p:Person {name: 'Alice'})-[:FRIENDS_WITH*2]-(fof)
WHERE fof <> p
RETURN DISTINCT fof.name

// Shortest path
MATCH path = shortestPath(
    (a:Person {name: 'Alice'})-[*..6]-(b:Person {name: 'Charlie'})
)
RETURN path
\end{minted}
\end{example}

\begin{remark}
Graph databases are ideal for social networks, fraud detection,
recommendation systems, and knowledge graphs. They outperform relational
databases when queries involve deep recursive joins.
\end{remark}

% ------------------------------------------------------------
\section{Comparison and selection criteria}
% ------------------------------------------------------------

\begin{center}
\begin{tabular}{lcccc}
\hline
& \textbf{Key-value} & \textbf{Document} & \textbf{Column} & \textbf{Graph} \\
\hline
Example     & Redis    & MongoDB    & Cassandra & Neo4j \\
Schema      & None     & Flexible   & Semi-fixed & Flexible \\
Queries     & By key   & Rich       & By partition & Traversals \\
Scalability & Very high & High      & Very high & Moderate \\
Typical use & Cache    & CMS, catalog & Time series & Social network \\
\hline
\end{tabular}
\end{center}

\begin{proposition}[Selection criteria]
The choice of NoSQL database type depends on:
\begin{enumerate}
  \item the \textbf{structure} of the data (fixed vs.\ flexible schema),
  \item \textbf{access patterns} (by key, complex queries, traversals),
  \item \textbf{consistency} vs.\ \textbf{availability} requirements,
  \item the \textbf{volume} and \textbf{velocity} of the data.
\end{enumerate}
\end{proposition}

% ------------------------------------------------------------
\section{Key formulas}
% ------------------------------------------------------------

\begin{keyformulas}
\begin{itemize}
  \item \textbf{CAP}: Consistency + Availability + Partition tolerance --- at most 2 out of 3
  \item \textbf{Key-value}: $O(1)$ access by key, no queries on values
  \item \textbf{Document}: flexible schema, secondary indexes, aggregation
  \item \textbf{Column-family}: query-driven modeling, massive distributed writes
  \item \textbf{Graph}: traversals in $O(k)$ where $k$ = number of edges traversed
\end{itemize}
\end{keyformulas}

% ------------------------------------------------------------
\section{Exercises}
% ------------------------------------------------------------

\begin{exercise}
For each of the following use cases, recommend a NoSQL database type and
justify: (a) web session cache, (b) product catalog with variable
attributes, (c) IoT sensor tracking at 100,000 events/s, (d) community
detection in a social network.
\end{exercise}

\begin{exercise}
Write a Python script that inserts 10,000 documents into MongoDB, creates
an index on a field, then compares the query time with and without the
index.
\end{exercise}

\begin{exercise}
Model a movie recommendation system in Neo4j. Create nodes (users, movies,
genres) and relationships (WATCHED, RATED, HAS\_GENRE). Write a Cypher
query to recommend movies to a user based on their friends' preferences.
\end{exercise}

\begin{exercise}
Design a Cassandra schema for a messaging application (messages between
users). The primary query is ``display the last $N$ messages in a
conversation.'' Justify the choice of partition key and clustering columns.
\end{exercise}

\begin{exercise}
Discuss the implications of the CAP theorem for a banking application
distributed across multiple continents. What trade-off would you
recommend? Why?
\end{exercise}
